\subsection{RQ1: $\phi$ as a static analysis}
\label{sec:rq_analysis}

\paragraph{Design.} An abstraction that is evaluated only by what it does for a
downstream classifier can hide errors in its own labels. We therefore score
$\phi$'s callee classification against annotated ground truth. This version of
the validation differs from the one in the first submission of this work in
three respects, each responding to a concern that a validation of this kind is
only as good as its annotation.

\begin{itemize}
    \item \textbf{The experimental corpora.} Callee names and one
          example call each are drawn from the \emph{experimental samples} of
          BigVul and DiverseVul, not only from Devign; Devign (not used in the
          detection experiments) is kept so that the new labels can be compared
          with the first submission's. The sample is stratified
          over $\phi$'s predicted classes (up to 14 names per class, half from
          the most frequent names and half at random) plus 100 names drawn
          from callees that $\phi$ left unlabelled, so that recall is
          estimable. The Devign gold set of the original submission is
          annotated again.
    \item \textbf{Independent, blind annotation.} Each name is labelled by two
          annotators who see only the name, its frequency and one example call;
          they never see $\phi$'s prediction, each other's labels, or the
          analysis code. The criteria are written down in advance
          (\texttt{experiments/annotations/PROTOCOL.md}). The two annotators
          are language-model annotators of different capability, run as
          separate sessions. No human annotated these sheets.
    \item \textbf{Agreement and adjudication.} We report raw agreement and
          Cohen's $\kappa$~\cite{cohen1960coefficient} between the two annotators; every disagreement is
          resolved by a third blind pass that sees the two candidate labels in
          random order. The consensus is the ground truth for scoring. We also
          score against each annotator alone to show sensitivity.
\end{itemize}

\paragraph{Who annotated the original gold set.} The 212 Devign labels of the
first submission were assigned by a single annotator from written criteria. The sheet
that annotator worked from listed $\phi$'s predicted class next to each name,
so the annotation was not blind to the analysis under test, and the labels are
stored in a script in which every name not listed defaults to ``none''. That
is a weaker protocol than the blind one above.
Their agreement with the new consensus is $\kappa_C=\numdevignOrigKappa$,
so the original labels were largely sound but are not used for the headline
scores here.

% Generated by experiments/emit_tables.py -- DO NOT EDIT BY HAND.
% Every value is read from experiments/results/*.json.

\begin{table}[t]
\centering
\caption{RQ1: $\phi$ as a static analysis, scored against the consensus of two
blind annotators and a third pass on disagreements, pooled over the three
corpora. Brackets are 95\% Wilson intervals; tp/fp/fn are counts of distinct
callee names. The final column group reports the question $\phi_8$ actually
answers (memory or string operation vs. not).}
\label{tab:rq1_analysis}
\begin{tabular}{lrrr}
\toprule
Class & Precision & Recall & tp/fp/fn \\
\midrule
\texttt{MEMORY\_ALLOC} & 0.88 $[0.75,0.94]$ & 0.78 $[0.65,0.87]$ & 42/6/12 \\
\texttt{MEMORY\_REALLOC} & 0.97 $[0.85,0.99]$ & 0.94 $[0.82,0.98]$ & 34/1/2 \\
\texttt{MEMORY\_FREE} & 0.85 $[0.73,0.93]$ & 0.95 $[0.85,0.99]$ & 41/7/2 \\
\texttt{MEMORY\_COPY} & 0.63 $[0.46,0.77]$ & 0.88 $[0.70,0.96]$ & 22/13/3 \\
\texttt{MEMORY\_SET} & 1.00 $[0.83,1.00]$ & 0.95 $[0.76,0.99]$ & 19/0/1 \\
\texttt{STRING\_COPY} & 1.00 $[0.90,1.00]$ & 0.82 $[0.68,0.91]$ & 33/0/7 \\
\texttt{STRING\_CONCAT} & 0.97 $[0.83,0.99]$ & 0.93 $[0.79,0.98]$ & 28/1/2 \\
\texttt{STRING\_FORMAT} & 0.86 $[0.71,0.94]$ & 0.88 $[0.73,0.95]$ & 30/5/4 \\
\texttt{STRING\_LENGTH} & 0.97 $[0.83,0.99]$ & 0.97 $[0.83,0.99]$ & 29/1/1 \\
\texttt{IO\_CALL} & 0.86 $[0.74,0.93]$ & 0.69 $[0.57,0.79]$ & 43/7/19 \\
\midrule
Macro & 0.90 & 0.88 & \\
Memory operation (any of the 9) & 0.94 & 0.94 & 293/19/19 \\
\bottomrule
\end{tabular}
\end{table}

\begin{table}[t]
\centering
\caption{RQ1: annotation agreement and per-corpus scores. Agreement and
Cohen's $\kappa$ are between the two blind annotators before adjudication;
precision and recall are macro-averaged over classes against the consensus; the
last column is precision / recall for ``memory or string operation''.}
\label{tab:rq1_agreement}
\resizebox{\linewidth}{!}{%
\begin{tabular}{lrrrrrr}
\toprule
Corpus & Names & Agreement & $\kappa$ & Macro P & Macro R & Memory-op P / R \\
\midrule
Devign & 212 & 0.92 & 0.89 & 0.90 & 0.87 & 0.90 / 0.93 \\
BigVul & 227 & 0.94 & 0.93 & 0.90 & 0.87 & 0.95 / 0.95 \\
DiverseVul & 236 & 0.92 & 0.90 & 0.89 & 0.89 & 0.96 / 0.94 \\
\bottomrule
\end{tabular}}
\end{table}


\paragraph{Agreement.} Table~\ref{tab:rq1_agreement} reports agreement per corpus.
Before adjudication the two blind annotators agree on \numpooledAgree{} of the
\numpooledNNames{} names (pooled Cohen's kappa $\kappa_C=\numpooledKappa$; per
corpus, raw agreement and $\kappa_C$: Devign \numdevignAgree{} and
\numdevignKappa, BigVul \numbigvulAgree{} and \numbigvulKappa, DiverseVul
\numdiversevulAgree{} and \numdiversevulKappa), conventionally read as almost
perfect agreement. \numdevignAdjudicated, \numbigvulAdjudicated{} and
\numdiversevulAdjudicated{} disagreements respectively went to the third pass; in
\numdevignUnresolved{} Devign case(s) the third annotator agreed with neither of
the first two and its label was used.

\paragraph{Accuracy of $\phi$.} Against the consensus, macro-averaged precision
is \numpooledMacroP{} and recall \numpooledMacroR{} over the ten callee classes
(Table~\ref{tab:rq1_analysis}); these are lower than the $0.976$ reported for
the original single-annotator labels, which we now regard as optimistic. The
weakest class is \texttt{MEMORY\_COPY}, with precision \numcopyPrecision{}. Two
caveats keep these figures honest. First, several classes have few positives
(the smallest has \numminClassPositives), so their intervals are wide;
Table~\ref{tab:rq1_analysis} shows Wilson intervals~\cite{wilson1927probable} rather than only the point
values. Second, the sample is stratified, so the numbers are per-name
estimates, not a population-weighted accuracy.

\paragraph{What $\phi_8$ actually needs.} At $|\mathcal{T}|=8$ every memory and
string class collapses into a single \texttt{MEMORY\_ACCESS} type, so confusing
an allocation with a copy costs nothing. The question that matters is whether a
callee is a memory or string operation at all. On that question precision is
\numpooledMemopP{} and recall \numpooledMemopR{} pooled over the three corpora
(\numbigvulMemopP{}/\numbigvulMemopR{} on BigVul,
\numdiversevulMemopP{}/\numdiversevulMemopR{} on DiverseVul).

\paragraph{Residual disagreements.} The informative ones are structural, not
noise. QEMU's \texttt{tcg\_temp\_free*} releases an intermediate-representation
temporary rather than heap memory, so a release rule fires on it. Wrappers that
allocate \emph{and} copy or format (\texttt{g\_strdup},
\texttt{g\_strdup\_printf}, \texttt{string\_sprintf}) cannot be given a single
class by a single-label taxonomy. Project-specific I/O and teardown helpers
(\texttt{bdrv\_pread}, \texttt{ReadBlobMSBLong}, \texttt{RelinquishMagickMemory})
are missed because they match no allow-list entry and no whole-identifier-
component rule. These are limitations of a name-based analysis with no type
information; they bound what $\phi$ can claim and are one reason we do not
expect it to beat the parser's own alphabet.

\paragraph{What the validation changed in the design.} Matching whole identifier
components instead of substrings removes writes and predicates mislabelled as
releases and allocations (\texttt{\_\_put\_user}, \texttt{skb\_put},
\texttt{gen\_new\_label}, \texttt{float64\_is\_zero}); separating stream I/O
from buffer formatting keeps 655 \texttt{printf}/\texttt{fprintf} call sites in
one corpus out of the memory class; and restricting the copy rule to explicit
copy primitives raises its precision at the cost of recall. These three changes
predate the present evaluation and are unchanged.
